\section{Introdução}
\label{sec:introducao}

Sistemas distribuídos frequentemente recorrem a um processo coordenador para
organizar ações que dependem de uma visão única ou sequenciamento centralizado.
Como processos estão sujeitos a falhas arbitrárias e não anunciadas, essa
coordenação deve ser combinada com um detector de falhas capaz de identificar a
indisponibilidade do líder e acionar um protocolo de eleição para restabelecer
o serviço.

Este trabalho implementa o algoritmo clássico de eleição Bully (Valentão) e um detector
de falhas periódico baseado em \textit{heartbeat}~\cite{tanenbaum2023}, avaliados sobre contêineres
Docker independentes comunicando-se via gRPC com \textit{deadlines} explícitos.
O nó ativo de maior identificador detém a prioridade de liderança. Cada seguidor
verifica periodicamente o coordenador; ao constatar duas falhas consecutivas,
inicia uma eleição e notifica os demais participantes.

Além de validar as propriedades de correção no cenário base e examinar o impacto
do intervalo de \textit{heartbeat} no tempo de convergência, o trabalho investiga
o comportamento do sistema ao escalar o cluster para $N \in \{3, 5, 7, 10\}$ nós.
Especificamente, avalia-se:
\begin{enumerate}
    \item o crescimento do número de mensagens trocadas na eleição em função de
    $N$, confrontando a complexidade teórica $\mathcal{O}(N^2)$ com os dados
    empíricos coletados;
    \item a estabilidade do tempo de convergência algorítmica ao variar o tamanho
    da topologia;
    \item a resiliência do protocolo diante de cenários adversos, incluindo a
    queda simultânea dos dois nós de maior prioridade e a interrupção de nós
    subordinados intermediários.
\end{enumerate}
